%%%%%%%%%%%%%%%%
%%% lev 1 %%%
%%%%%%%%%%%%%%%%

\Chapter{1}

\Verse{1}
{
	\hP{וַיִּקְרָא אֶל מֹשֶׁה וַיְדַבֵּר יְהֹוָה אֵלָיו מֵאֹהֶל מוֹעֵד לֵאמֹר׃}
}
{
	\eP{
		And He called to Moses, and YHWH spoke to him from the
		Tent of Meeting, saying,
	}
}
{
	No turning back now.
	Strap in folks.

	Here we go.

	And he called\footnote{Hey that's the name of the book.} to Drew from the Tent of Meeting.

	The text is explicit: YHWH is in the tent.
}

\Verse{2}
{
	\hP{דַּבֵּר אֶל בְּנֵי יִשְׂרָאֵל וְאָמַרְתָּ אֲלֵהֶם אָדָם כִּי יַקְרִיב מִכֶּם קרְבָּן לַיהֹוָה מִן הַבְּהֵמָה מִן הַבָּקָר וּמִן הַצֹּאן תַּקְרִיבוּ אֶת קרְבַּנְכֶם׃}
}
{
	\eP{
		“Speak to the children of Israel. And you shall say to
		them: A human from you who will make an offering to
		YHWH—you shall make your offering from the domestic
		animals: from the herd and from the flock.
	}
}
{
	So YHWH tells Drew to talk to the people.

	And he's supposed to say this:

	You humans, if you want to give me something, it has to be house animals.

	So basically cows and sheep and so on.

	He doesn't want any wild animals as presents.

	Noted.

%	\footnote{
%		A human. The laws of sacrifice begin by referring to the offerer by the
%		species term “a human” (Hebrew ’dm). Rashi gives a midrashic explanation
%		for this. But we should note that the text next refers to the offerer as
%		“a person” (Hebrew nepe, Lev 2:1) and then as “a man” (Hebrew ’î, Lev
%		7:8), and this is the same order in which these terms are used to refer
%		to human beings in the book of Genesis (’dm in Gen 1:26; nepe in 2:7; ’î
%		in 2:23). This makes a link between the books of the Torah, and it makes
%		a link between God’s creation of humans on the one hand and humans’
%		sacrifices to their creator on the other. Creation and sacrifice are
%		linked. Creation produces life. Sacrifice takes life. And so the word
%		order here subtly reminds us that, when we take animals’ lives in order
%		to support our own, we should remember that both their lives and ours
%		are part of a common creation, and that both are treated in the Torah as
%		sacred.
%	}
}

\Verse{3}
{
	\hP{אִם עֹלָה קרְבָּנוֹ מִן הַבָּקָר זָכָר תָּמִים יַקְרִיבֶנּוּ אֶל פֶּתַח אֹהֶל מוֹעֵד יַקְרִיב אֹתוֹ לִרְצֹנוֹ לִפְנֵי יְהֹוָה׃}
}
{
	\eP{
		“If his offering is a burnt offering from the herd, he
		shall make it an unblemished male. He shall bring it
		forward to the entrance of the Tent of Meeting for his
		acceptance in front of YHWH.
	}
}
{
	If you want to burn the present, it has to be a perfect dude.

	No burning crippled animals or girls as a gift.

	You have to bring it to the door of the Tent of Meeting.

	Remember: YHWH's in the tent.

%	\footnote{
%		to the entrance of the Tent of Meeting. From this time on, there is only
%		one place on earth where an Israelite may perform a sacrifice: at the
%		altar at the entrance of the Tent of Meeting. This applies here to a
%		burnt offering. Later in this section we learn that it also applies to
%		the peace offering (3:2), the sin offering (4:3), the guilt offering
%		(7:2), and the grain offering (6:7–9). This applies to the period of
%		Israel’s travels through the wilderness. Then the tent is set up at
%		Shiloh (Josh 18:1), then at Gibeah (2 Chr 1:3), and then finally it is
%		brought to the Temple that King Solomon builds in Jerusalem (1 Kings
%		8:4; 1 Chr 5:5). One exception: the prophet Elijah makes a sacrifice at
%		Mount Carmel in a duel with the priests of Baal; but it is presented as
%		an extraordinary occurrence, led by a prophet, involving a miracle, the
%		last miracle in the Tanak that is said to be performed in front of the
%		entire people of Israel (1 Kings 18). All other sacrifice is to take
%		place at the Tent of Meeting. This centralization of worship in Israel’s
%		religion is crucially important. It is a powerful corollary of Israelite
%		monotheism: one God, one altar.
%	}
}

\Verse{4}
{
	\hP{וְסָמַךְ יָדוֹ עַל רֹאשׁ הָעֹלָה וְנִרְצָה לוֹ לְכַפֵּר עָלָיו׃}
}
{
	\eP{
		And he shall lay his hand on the head of the burnt
		offering, and it will be accepted for him, to atone for
		him.
	}
}
{
	Put your hand on its head.
	That makes it acceptable.

	For forgiveness.
}

\Verse{5}
{
	\hP{וְשָׁחַט אֶת בֶּן הַבָּקָר לִפְנֵי יְהֹוָה וְהִקְרִיבוּ בְּנֵי אַהֲרֹן הַכֹּהֲנִים אֶת הַדָּם וְזָרְקוּ אֶת הַדָּם עַל הַמִּזְבֵּחַ סָבִיב אֲשֶׁר פֶּתַח אֹהֶל מוֹעֵד׃}
}
{
	\eP{
		And he shall slaughter the herd animal in front of YHWH,
		and the sons of Aaron, the priests, shall bring the blood
		forward and fling the blood on the altar, all around,
		which is at the entrance of the Tent of Meeting.
	}
}
{
	Then you kill it.

	In front of YHWH, who's in the tent.

	Then the Aaronsons fling the blood all around.

	So far so good.

%	\footnote{
%		he shall slaughter … and the priests shall bring the blood forward. The
%		individual slaughters the animal; the priests handle the blood. This is
%		consistent with the conception found elsewhere in the Torah that blood
%		is sacred. It can contaminate, it can protect, and it “is the life.”
%	}
}

\Verse{6}
{
	\hP{וְהִפְשִׁיט אֶת הָעֹלָה וְנִתַּח אֹתָהּ לִנְתָחֶיהָ׃}
}
{
	\eP{
		And he shall flay the burnt offering and cut it into its
		parts.
	}
}
{
	Then you cut it up.
}

\Verse{7}
{
	\hP{וְנָתְנוּ בְּנֵי אַהֲרֹן הַכֹּהֵן אֵשׁ עַל הַמִּזְבֵּחַ וְעָרְכוּ עֵצִים עַל הָאֵשׁ׃}
}
{
	\eP{
		And the sons of Aaron, the priest, shall put fire on the
		altar and arrange wood on the fire.
	}
}
{
	1. Fire on altar.
	2. Wood on fire.

	We're considering switching the order of those.

%	\footnote{
%		the sons of Aaron, the priest. Leviticus’s laws, as well as its
%		narrative, are fundamentally concerned with priesthood, and thus with
%		leadership, authority—human authority, but sanctioned by
%		God—transforming the Israelites from a ragtag mass of ex-slaves into an
%		organized society. The priests are installed and anointed, their
%		sacrificial function is institutionalized, and they are commissioned to
%		teach (10:10–11), a commission that implies authority. The installation
%		of the priesthood—Aaron and his sons—is divinely stamped by the only
%		appearance of the divine glory (the kbôd YHWH) in Leviticus, at the end
%		of eight days of inaugural ceremonies of the priesthood and the
%		sacrificial system. And it is emphasized that the priestly laws applying
%		to Aaron and his sons are eternal (10:9,15). The consecration of the
%		priesthood had been announced in summary in Exodus 40. Now it is
%		narrated. The priests’ ritual and ethical activity is tied to sacrifice,
%		and so the narrative’s first concern is the new priesthood’s first
%		sacrificial offering. But sacrifice cannot be performed in any way one
%		wishes. One must perform it in a prescribed manner. And so first there
%		are seven chapters of prescription. This is consistent with a theme of
%		prescription that pervades the book. There is a way to behave: in
%		morality, in ethics, and in ritual.
%	}
}

\Verse{8}
{
	\hP{וְעָרְכוּ בְּנֵי אַהֲרֹן הַכֹּהֲנִים אֵת הַנְּתָחִים אֶת הָרֹאשׁ וְאֶת הַפָּדֶר עַל הָעֵצִים אֲשֶׁר עַל הָאֵשׁ אֲשֶׁר עַל הַמִּזְבֵּחַ׃}
}
{
	\eP{
		And the sons of Aaron, the priests, shall arrange the
		parts, the head, and the suet on the wood that is on the
		fire that is on the altar.
	}
}
{
	Then the Aaronsons, the Cohens, will organize the head and stuff on the wood that's on the fire that's on the altar.
}

\Verse{9}
{
	\hP{וְקִרְבּוֹ וּכְרָעָיו יִרְחַץ בַּמָּיִם וְהִקְטִיר הַכֹּהֵן אֶת הַכֹּל הַמִּזְבֵּחָה עֹלָה אִשֵּׁה רֵיחַ נִיחוֹחַ לַיהֹוָה׃}
}
{
	\eP{
		And he shall wash its innards and its legs with water. And
		the priest shall burn it all to smoke at the altar, a
		burnt offering, an offering by fire of a pleasant smell to
		YHWH.
	}
}
{
	Wash the guts and legs.
	Then once it's clean they burn it up.

	That's a burnt offering.

	By fire.

	YHWH likes the smell.
}

\Verse{10}
{
	\hP{וְאִם מִן הַצֹּאן קרְבָּנוֹ מִן הַכְּשָׂבִים אוֹ מִן הָעִזִּים לְעֹלָה זָכָר תָּמִים יַקְרִיבֶנּוּ׃}
}
{
	\eP{
		“And if his offering is from the flock—from the sheep or
		from the goats—for a burnt offering, he shall make it an
		unblemished male.
	}
}
{
	If you're burning sheep or goats, it has to be a perfect one. No cripples or girls.
}

\Verse{11}
{
	\hP{וְשָׁחַט אֹתוֹ עַל יֶרֶךְ הַמִּזְבֵּחַ צָפֹנָה לִפְנֵי יְהֹוָה וְזָרְקוּ בְּנֵי אַהֲרֹן הַכֹּהֲנִים אֶת דָּמוֹ עַל הַמִּזְבֵּחַ סָבִיב׃}
}
{
	\eP{
		And he shall slaughter it on the northward side of the
		altar in front of YHWH. And the sons of Aaron, the
		priests, shall fling its blood on the altar, all around.
	}
}
{
	Kill it on the north side in front of YHWH.
	Then they fling the blood all around again.
}

\Verse{12}
{
	\hP{וְנִתַּח אֹתוֹ לִנְתָחָיו וְאֶת רֹאשׁוֹ וְאֶת פִּדְרוֹ וְעָרַךְ הַכֹּהֵן אֹתָם עַל הָעֵצִים אֲשֶׁר עַל הָאֵשׁ אֲשֶׁר עַל הַמִּזְבֵּחַ׃}
}
{
	\eP{
		And he shall cut it into its parts and its head and its
		suet. And the priest shall arrange them on the wood that
		is on the fire that is on the altar.
	}
}
{
	Cut it up, arrange the head and all, on the wood that's on the fire that's on the altar.
}

\Verse{13}
{
	\hP{וְהַקֶּרֶב וְהַכְּרָעַיִם יִרְחַץ בַּמָּיִם וְהִקְרִיב הַכֹּהֵן אֶת הַכֹּל וְהִקְטִיר הַמִּזְבֵּחָה עֹלָה הוּא אִשֵּׁה רֵיחַ נִיחֹחַ לַיהֹוָה׃}
}
{
	\eP{
		And he shall wash the innards and the legs with water. And
		the priest shall bring it all forward and burn it to smoke
		at the altar. It is a burnt offering, an offering by fire
		of a pleasant smell to YHWH.
	}
}
{
	Wash the guts and legs.
	Then once it's clean they burn it up.

	That's a burnt offering.

	By fire.

	YHWH likes the smell.
}

\Verse{14}
{
	\hP{וְאִם מִן הָעוֹף עֹלָה קרְבָּנוֹ לַיהֹוָה וְהִקְרִיב מִן הַתֹּרִים אוֹ מִן בְּנֵי הַיּוֹנָה אֶת קרְבָּנוֹ׃}
}
{
	\eP{
		“And if his offering to YHWH is a burnt offering from
		birds, then he shall make his offering from turtledoves or
		from pigeons.
	}
}
{
	If you want to give YHWH some birds as a burning present, it has to be either:
	1. Turtledoves
	Or
	2. Pigeons
}

\Verse{15}
{
	\hP{וְהִקְרִיבוֹ הַכֹּהֵן אֶל הַמִּזְבֵּחַ וּמָלַק אֶת רֹאשׁוֹ וְהִקְטִיר הַמִּזְבֵּחָה וְנִמְצָה דָמוֹ עַל קִיר הַמִּזְבֵּחַ׃}
}
{
	\eP{
		And the priest shall bring it forward to the altar and
		wring its head and burn it to smoke at the altar, and its
		blood shall be drained on the wall of the altar.
	}
}
{
	The priest chokes the bird and burns it up.
	All its blood goes on the wall of the altar.
}

\Verse{16}
{
	\hP{וְהֵסִיר אֶת מֻרְאָתוֹ בְּנֹצָתָהּ וְהִשְׁלִיךְ אֹתָהּ אֵצֶל הַמִּזְבֵּחַ קֵדְמָה אֶל מְקוֹם הַדָּשֶׁן׃}
}
{
	\eP{
		And he shall take away its crop with its feathers and
		throw it beside the altar eastward to the place of the
		ashes.
	}
}
{
	Bird remains go in the big pile of ashes to the left.
	If you're looking at YHWH, turn left. The bird leftovers go there.
}

\Verse{17}
{
	\hP{וְשִׁסַּע אֹתוֹ בִכְנָפָיו לֹא יַבְדִּיל וְהִקְטִיר אֹתוֹ הַכֹּהֵן הַמִּזְבֵּחָה עַל הָעֵצִים אֲשֶׁר עַל הָאֵשׁ עֹלָה הוּא אִשֵּׁה רֵיחַ נִיחֹחַ לַיהֹוָה׃}
}
{
	\eP{
		And he shall tear it open by its wings—he shall not divide
		it—and the priest shall burn it to smoke at the altar on
		the wood that is on the fire. It is a burnt offering, an
		offering by fire of a pleasant smell to YHWH.
	}
}
{
	Don't cut it in half. Remember Abraham when he cut those animals up? He didn't do the birds right? Exactly.

	So just tear it open. The wings are like a pull here to open thing.

	Then you do the thing again, and YHWH likes the smell.

	Simple enough right?

	End of chapter.

%	\footnote{
%		he shall not divide it. One cuts up herd and flock animals during a
%		sacrifice, but one does not divide birds this way. This summons to mind
%		the original sacrifice that sealed the covenant between God and Abraham,
%		in which Abraham cuts all the animals in half but leaves the birds whole
%		(Gen 15:10), and a fire representing God’s presence passes between the
%		parts of the animals. Thus Israel’s manner of sacrifice now memorializes
%		the fact that God bound Himself to Israel through a ceremony of animal
%		sacrifice.
%	}
}
